\section{Filtered subrepresentations and character blocks at the two ramified primes}

% BEGIN H80M1
\paragraph{H80M1. The Fontaine subrepresentations at the distinguished primes.}
Work on the original hub of F78D2 at a twisted index with
$\ell=q$ or $d$ and $\mathrm{HG}$. These primes differ from the two
type primes. The good-point identity gives a point of $Z_i$ with
its continuous complement and the canonical family is good. Thus
S80J3 gives crystallinity of $V_x=R_x/\chi_1$ at every point, on
every coefficient extension of every $w\mid\ell$. Fix one such
extension $L'$ and write $K=E_w$. The dual matrices in the source
act on $V_x^\vee(1)$. These are crystalline by duality and the
cyclotomic period. For any representation $U$ in
$\mathrm{IsSubOfSums}$ there is a specific injective equivariant
map $U\to\bigoplus_{j=1}^s V_{x_j}^\vee(1)$. In particular
one can give $U$ the subspace topology via this map. The image is
a closed, stable finite-dimensional subspace, so its action is
continuous for its finite-dimensional topology.

Over the tied period rings the crystalline category is closed under
subobjects and quotients. One can see the dimension assertion
directly: invariants are left exact, and each term has crystalline
dimension at most its representation dimension over unramified
scalars, counting coefficients. For a short exact sequence with
crystalline middle term these two upper bounds for subobject and
quotient sum to the attained middle dimension. Hence both are
attained and the sequence on periods is exact. On de Rham spaces
the filtration sequence is strict. Indeed the subobject filtration is
induced, since tensoring the underlying vector-space injection and
its linear splitting with each $\mathrm{Fil}^zB_{\rm dR}$ gives
the intersection formula. This suffices here for an injection on
graded pieces. The standard strict de Rham comparison also gives
the quotient filtration. These statements transport through A72F2
to the specified coefficient embedding. Consequently every jump $a$
of $U$ occurs on some $V_{x_j}^\vee(1)$, and $-a-1$ is a jump
of $V_{x_j}$. If $s=0$, then $U=0$ and the implication is vacuous.

At each good $\pi$, the module $T_{2,\pi}$ is the Tate-twisted
dual of the rank complement in $V_\pi$. It is crystalline by S80J3.
By S80J1 its multiplicity at $a$ equals one precisely when
$-a-1$ is a jump of $V_\pi$ outside $[0,t-1]$, and equals zero
otherwise. This proves both assertions of
$\mathrm{FontaineSubOfSumsAtDistinguishedPrimes}$ on the original hub.
% END H80M1

% BEGIN H80M2
\paragraph{H80M2. Every row input to the ramified local lifting estimate.}
Write $b=L.\mathrm{sc}.b$, $S=S_c$, $\eta=\eta_{\rm Sz}$, and set
$H=\lfloor2b\rfloor$, $S_0=\lfloor S\rfloor$,
$A=\lfloor\eta\rfloor+t$. R79Q3 gives at each $\tau:K\to L'$
the integral cloud and low sets $J_c,J_l$ with $|J_l|=m$.
Its integer $s$ satisfies $\eta+t+1<s\le b+\eta+t+1$.
If the cloud is in $[s,s+S]$, duality places it in
$[-(s+1)-S_0,-(s+1)]$. Use the negative deep interval with
endpoint $s+1$. If the cloud is in $[-s-S,-s]$, use the positive
deep interval $[s-1,s-1+S_0]$. Either positive endpoint exceeds
$A$. The dual lows lie in
$[-\eta-(t+1)/2,\eta-(t+1)/2]$ by S80J2, and the dual string
in $[-t,-1]$. Their integer points all belong to $[-A,A]$.
The parameter estimate $S+\eta+t+2\le b$ puts the endpoint
$s+1+S_0$ at most $2b$. Thus all these integral jumps are also
in $[-H,H]$. By H80M1 this gives the deep-or-central and absolute
bounds for every subrepresentation allowed in the source.

For $|a|\le A$ the cloud is absent, and the string contributes
nothing to $T_{2,\pi}$. The remaining possible values have
$-a-1\in J_l$, each with multiplicity at most one. The involution
$a\mapsto-a-1$ injects them into the $m$-element list. Thus their
total graded dimension is at most $m$ at each $\tau$. In particular
the estimate counts dimensions even if the source lattice has
arbitrary rank.

For the native row map the source representation is a subobject of
such a sum, the target is given by an integral matrix-lattice
isomorphism with the actual $T_{2,\pi}$, and the map is onto in
the specified sense. The target has rank $r$ and the same filtered
multiplicities. The last alternative in the local lifting theorem is
automatic because $[K:K_0]=2$ by R79Q2. Its good level-one
condition follows from $K\subset K_0(\mu_\ell)$, the good
coefficient-extension hypothesis and oddness of $\ell$.
Thus, if the original local theorem $\mathrm{thmA}$ with constant
$C_A$ is supplied on these fields, the following rational row
length works for each declaration, all its good coefficient extensions
and every allowed row map, where $d_w=[K:\mathbb Q_\ell]$:
\[
 \left\lceil\max(C_A,0)
      \left((d_w+1)r\bigl(4b/\ell+2S+2\log_\ell(2b)+1+\eta+t\bigr)
                   +3d_wm(2b+1)\right)\right\rceil .
\]
Indeed insert $e=2$, $n=[K_0:\mathbb Q_\ell]$,
$(S_0,A,m,H)$ into $\mathrm{boundA}$, use $2n=d_w$, and increase
each floor to its real upper bound. The logarithm is increasing and
$H\ge1$. The ceiling provides a rational majorant less than the
displayed real expression before rounding plus one. This gives the
original $\mathrm{RowAt}$ inequality on the tested integral maps
with their own exponent and index witnesses from $\mathrm{thmA}$.
% END H80M2

% BEGIN H80M3
\paragraph{H80M3. An explicit joint tame exponent before the character estimate.}
Keep $H=\lfloor2b\rfloor$ and put
$N=4n(r+t)$ with $n$ the global rank parameter, and
$P=\lfloor\log_\ell(2H)\rfloor$. At each $w$ and every good
level-one coefficient field, take any crystalline $Y$ whose
labelled jumps are contained in those of $V_\pi$ for a point
$\pi$ at the given index. Every label has at most $n_V=r+t$
distinct jumps in $[-H,H-1]$. The upper endpoint follows from
R79Q3 and $S+\eta+t+2\le b$: on the positive cloud $z+1\le2b$,
on the negative cloud $-z\le2b$, and the lows and string obey
both bounds. The coefficient field splits $K$ and it has $d_w$
labels. The total jump set of $Y$ is contained in their union:
$L'\otimes_{\mathbb Q_\ell}K$ splits into the labelled factors,
so a nonzero graded piece has a nonzero component. The union
has at most $d_wn_V\le N$ elements, using $d_w\le2f=4n$.

The good layer has $s_{\rm loc}=1$. R79Q2 gives
$[K_0(\mu_\ell):K]=(\ell-1)/2$, of $\ell$-valuation zero, hence
$g_K=0$. From the Eisenstein cyclotomic degrees,
$[K_0(\mu_{\ell^j}):K_0(\mu_\ell)]=\ell^{j-1}$ for every $j\ge1$.
Consequently $\mathrm{vdegMs}(j)=j-1$. The inertia group in the
bridge is $I_K$ because its decomposition-group identification
preserves the valuation and its induced residue action.

Write $U$ for the total jump set of $Y$ and
$M=3+\min\{j\ge0:2H\le\ell^j(\ell-1)\}$. Then $M\le4+P$.
For $0\le i\le e<t\le H$, each $u\in U$ different from $i$
satisfies $0<|u-i|\le2H$, hence $v_\ell(u-i)\le P$. It follows that
$\sum_{u\ne i}(M+v_\ell(u-i))\le N(4+2P)$. In the source's
$\mathrm{cLL}$ sum the term for $i$ adds two to the maximum of one
and this quantity, and there are $e+1\le t$ terms, with an initial
degree contribution $M-1$. The final $\mathrm{cTB}$ corrections
at $s_{\rm loc}=1$, $g_K=0$ consist of
$(e+2)+v_\ell(e)+\lfloor e/(\ell-1)\rfloor+1+1+v_\ell(e+1)$.
With the native convention $v_\ell(0)=0$, these corrections
are at most $4t+3$ using $v_\ell(e)\le e$ and
$v_\ell(e+1)\le e+1$. Thus every $\mathrm{cTBY}(Y,H,H,e)$ is at most
\[
               c_B=(3+P)+t\{3+N(4+2P)\}+4t+3 .
\]
Together with $K\subset K_0(\mu_\ell)$ these are all the assertions
of $\mathrm{JointConditionTameBound}(H,H,c_B)$. They hold for every
$\pi$ at the index. By the numerical proof of C78B3 this same
integer obeys $c_B\le120tn(r+t)\log b=60tf(r+t)\log b$.
Its formula and this comparison depend only on the original member
and $\ell$.
% END H80M3

% BEGIN H80M4
\paragraph{H80M4. All three rank-one cohomology clauses at a ramified prime.}
For every $0\le e<t$ let $V=L'(-e)$, let $T$ be its specified
integral line, and put $A_e=V/T$. Every matrix-lattice realization
of the inverse cyclotomic line is integrally isomorphic to this one.
The mod-$\ell$ cyclotomic image on $G_K$ has order
$[K(\mu_\ell):K]=(\ell-1)/2$ by R79Q2. As $\ell>4tr+1$, the
exponent $e+1$ lies strictly between zero and this order.
Choose an element on which $\varepsilon_\ell^{e+1}-1$ is a unit.
It also acts without invariants on the divisible dual twist of $T$,
a direct sum of copies of $\mathbb Q_\ell/\mathbb Z_\ell(e+1)$
after restriction of coefficient scalars. Local Tate duality on
the original lattice gives $H^2(K,T)=0$ and rationally
$H^2(K,V)=0$.

The long exact sequence for $T\to V\to A_e$ now identifies
$H^1(K,A_e)$ with $H^1(K,V)$ modulo the image lattice of $H^1(K,T)$.
Here integral local cohomology is finitely generated, and inverting
$\ell$ recovers rational cohomology, so the image is a full
$O'$-lattice. The rational Euler characteristic gives
$\dim_{L'}H^1(K,V)=d_w+\mathbf1_{e=0}$, as the cyclotomic
character has infinite image. The Bloch--Kato finite subgroup
on $A_e$ is the image of the rational finite subspace, by the
definition with rational lifts and torsion coboundaries.
Consequently its quotient in $H^1(K,A_e)$ is a finite-dimensional
$L'$-space modulo a full lattice of rank at most $d_w+1$.
The part killed by $\ell^a$ therefore has length at most
$(d_w+1)a\,\mathrm{length}_{O'}(O'/\ell O')$ for every $a\ge0$.
The map taking a cocycle to its class identifies this torsion
submodule with the quotient of cocycles whose $\ell^a$-multiples
are finite by the finite cocycles. Its kernel is precisely the
finite cocycles and every class has a cocycle representative.
This is the original length clause, in particular at $a=c_B$.

For $e>0$ the jumps of $V$ equal $e$, so
$D_{\rm dR}(V)/\mathrm{Fil}^0=0$, and $H^0(K,V)=0$.
The Bloch--Kato dimension formula gives $H_f^1(K,V)=0$.
Thus each finite cocycle on $A_e$ is a coboundary.
For $e=0$ the same formula gives dimension one. The unramified
rational classes are a line, all finite: the unipotent extension
corresponding to such a homomorphism is an unramified representation,
hence crystalline by the invariant Witt basis argument of S80J4.
This proves equality of the rational finite and unramified lines.
An unramified homomorphism into $L'/O'$ lifts through this rational
line: lift its Frobenius value into $L'$ and multiply it by the
continuous projection $\widehat{\mathbb Z}\to\mathbb Z_\ell$.
The equality holds on all the unramified quotient by density.
The coboundaries of the trivial action are zero. This proves the
required equality of finite and inertia-trivial cocycles and
completes the second half of $\mathrm{TBAt}$, for every character
tested at the native ramified row.
% END H80M4

% BEGIN H80M5
\paragraph{H80M5. The remaining integral test and the numerical sum at the two primes.}
At the ramified places the first half of $\mathrm{TBAt}$ would follow
from $\mathrm{TBat}$ on the good level-one layers of this very Wach
functor. In fact for a tested crystalline lattice $Y'$ with
invariant projection $p$ and integral vector $e_0$ above one,
let $Y=\ker p$ with its intersection lattice. The integral sequence
onto the trivial line is exact because $e_0$ splits its underlying
modules. The rational kernel is crystalline and its jumps lie in
those of $Y'$, by strictness as in H80M1. H80M3 puts both in the
height interval and bounds $c_{\rm TBY}(Y,H,H,e)$ by $c_B$.
W79H5 gives their $\mathrm{HypW}$ records, in particular the
required restricted crystallinity for $Y'$. In $\mathrm{TBat}$
take the integral lift $e_0$, the kernel cocycle $g e_0-e_0$ and
the prescribed equivariant map from the kernel lattice into $A_e$.
The asserted Bloch--Kato finiteness after
$\ell^{c_{\rm TBY}}$ implies it after $\ell^{c_B}$ by multiplication
with the remaining nonnegative power. With H80M4 this gives the
original $\mathrm{admTB}$ condition if that local kernel-cocycle
estimate is supplied.

For clarity, sum these quantities in the all-row declaration of
R79Q4 under the row lifting estimate $\mathrm{thmA}$ and this
kernel-cocycle estimate. Put $u=\log b>1$, $K_A=\max(C_A,0)$,
$B=b/\ell+S+t+\eta+u$ and $Q_0=f^2t^2(r+t)$.
Since $\sum_w d_w\le2f$ and $\#\{w\}\le f$, H80M2 bounds the
sum of row lengths by
\[
 K_A f\!\left(3r(4b/\ell+2S+2\log_\ell(2b)+1+\eta+t)
                      +6m(2b+1)\right)+f .
\]
This is at most $13K_A f(rB+mb)+f$. Indeed
$\log_\ell(2b)\le2u$ since $\ell\ge3$ and $u>1$.
In comparing coefficients the terms with $b/\ell,S,u,\eta,t,mb$
then have coefficients at most $12,6,12,3,3,12$ respectively.
The constant $3r+6m\le9r$ is absorbed by the remaining
$10rt$, using $m\le r$, $t\ge1$.

The block error $t\sum_w(d_w+2)c_B$ is at most $240Q_0u$
by H80M3. The final valuation error $4n_Vv_\delta(i)$ is at most
$4(r+t)t$ by R79Q4. As $f\le Q_0$ and $(r+t)t\le Q_0$, all
the errors together are at most
\[
             13K_A f(rB+mb)+260Q_0u .
\]
On the native global lattice-index inequality for these data, the
remaining terms are the Selmer length, the positive twists
$1\le j<t$, the class-number term and the root-of-unity correction.
Killing bounds this last correction by $\ell^{-2}$, since R79Q4
gives $tr<[E(\mu_\ell):E]$. At $C_A=160$ the coefficient is
$13\cdot160=2080$. J78H2 then applies at its original polynomial
and constant. This calculation isolates the integral row lifting
at $160$, the kernel-cocycle estimate in the first half of
$\mathrm{TBAt}$, and application of the global index inequality on
the all-row data as remaining requirements of this argument.
% END H80M5
